% Pollis New Methods: two-scales realized volatility (Zhang, Mykland and Aït-Sahalia) — draft 1 (11 November 2026)
% Compile with: latexmk -pdf TwoScalesVolatility.tex && latexmk -c
% Shared look for the Pollis New Methods documents (same fonts and colours as the brochure and the workflow handouts).
\documentclass[10pt,a4paper]{article}

\usepackage[a4paper,left=18mm,right=18mm,top=15mm,bottom=17mm,footskip=8mm]{geometry}
\usepackage[T1]{fontenc}
\usepackage[default,semibold]{sourcesanspro}
\usepackage{textcomp}
\usepackage[italic]{mathastext}
\usepackage{amsmath}
\usepackage{microtype}
\usepackage[table]{xcolor}
\usepackage{tikz}
\usepackage[most]{tcolorbox}
\usepackage{fontawesome5}
\usepackage{tabularx,booktabs}
\usepackage{enumitem}
\usepackage{fancyhdr}
\usepackage[hidelinks]{hyperref}

% ---------------------------------------------------------------- colours (as in the brochure)
\definecolor{navy}{HTML}{1C3A5E}
\definecolor{pblue}{HTML}{3B78B5}
\definecolor{amber}{HTML}{F2B33D}
\definecolor{deepamber}{HTML}{D9902A}
\definecolor{pale}{HTML}{EEF3F9}
\definecolor{ink}{HTML}{2A2F36}
\definecolor{logonavy}{HTML}{173D63}
\definecolor{logodark}{HTML}{C98A1E}
\definecolor{logoamber}{HTML}{F2B446}
\definecolor{logocream}{HTML}{FEEFD2}
\definecolor{soft}{HTML}{5E6773}
\definecolor{agree}{HTML}{2F7D4F}
\definecolor{differ}{HTML}{B23B3B}
\color{ink}

\setlength{\parindent}{0pt}
\setlength{\parskip}{4pt}
\setlist{nosep,leftmargin=*,itemsep=2pt}
\renewcommand{\arraystretch}{1.15}

\pagestyle{fancy}
\fancyhf{}
\renewcommand{\headrulewidth}{0pt}
\fancyfoot[L]{\footnotesize\color{soft}\textbf{Pollis}\enspace New Methods\enspace·\enspace\nmtheme}
\fancyfoot[R]{\footnotesize\color{soft}\thepage}
\newcommand{\nmtheme}{}
\newcommand{\nmdate}{11 October 2026}

% The Pollis logo, redrawn so it prints sharply at any size
\newcommand{\pollislogo}[1]{%
  \begin{tikzpicture}[x=#1/512,y=-#1/512,baseline=(current bounding box.center)]
    \fill[logonavy] (0,0) rectangle (512,512);
    \fill[logodark] (36,422) -- (104,256) -- (172,422) -- cycle;
    \fill[logoamber] (104,422) -- (203,180) -- (302,422) -- cycle;
    \fill[logocream] (203,422) -- (339,89) -- (476,422) -- cycle;
  \end{tikzpicture}}

% Top of the first page: logo, name; then the kicker, title and a one-line promise
\newcommand{\nmhead}[3]{%
  \renewcommand{\nmtheme}{#1}%
  \thispagestyle{fancy}%
  \begin{minipage}[c]{13mm}\pollislogo{11mm}\end{minipage}%
  \begin{minipage}[c]{\dimexpr\linewidth-13mm\relax}
    {\color{navy}\Large\bfseries Pollis}\enspace{\color{soft}New Methods}\hfill
    {\color{deepamber}\small\bfseries\MakeUppercase{#1}}
  \end{minipage}\par\vspace{7mm}
  {\color{navy}\fontsize{22}{26}\selectfont\bfseries #2}\par\vspace{4pt}
  {\color{soft}\large #3}\par\vspace{5mm}}

\newcommand{\nmsection}[1]{\par\vspace{5pt}{\color{navy}\large\bfseries #1}\par\vspace{1pt}}
\newcommand{\nmsub}[1]{\par\vspace{2pt}{\color{deepamber}\bfseries #1}\par}
\newcommand{\gt}{\,{\color{deepamber}\bfseries ›}\,}
\newcommand{\menupath}[1]{{\bfseries\color{navy}#1}}
\newcommand{\fn}[1]{\texttt{\small #1}}
% Citations: \nmcite{key}{Breiman (2001)} links to the entry \nmref{key}{...} in the References
\newcommand{\nmcite}[2]{\hyperlink{ref:#1}{\color{pblue}#2}}
\newcommand{\nmref}[2]{\raisebox{1.1em}[0pt][0pt]{\hypertarget{ref:#1}{}}#2\par}
\newcommand{\nmdoi}[1]{\href{https://doi.org/#1}{\color{pblue}doi:#1}}
% Where a notebook is in Pollis: its title in Browse Methods, and its file
\newcommand{\browse}[2]{\menupath{Compose \gt New Methods \gt Browse Methods}: \emph{#1} (\fn{#2.ipynb})}

% Key facts at the top: source, missing in Julia, notebook, result
\newcommand{\factsheet}[4]{%
  \begin{tcolorbox}[enhanced, colback=pale, colframe=pale, boxrule=0pt, sharp corners,
    left=7pt, right=7pt, top=5pt, bottom=5pt, fontupper=\small]
  \begin{tabularx}{\linewidth}{@{}>{\leavevmode\bfseries\color{navy}}p{27mm}>{\leavevmode}X@{}}
  Tested against & #1\\
  Missing in Julia & #2\\
  Notebook & #3\\
  Result & #4
  \end{tabularx}
  \end{tcolorbox}}

% Boxes
\newtcolorbox{pbox}[2][]{enhanced, colback=pale, colframe=amber, boxrule=0pt,
  leftrule=3pt, sharp corners, left=7pt, right=7pt, top=6pt, bottom=6pt,
  fontupper=\small\raggedright, before upper={{\color{navy}\normalsize\bfseries #2}\par\smallskip}, #1}
\newtcolorbox{nbox}[2][]{enhanced, colback=navy, colframe=navy, boxrule=0pt,
  sharp corners, left=8pt, right=8pt, top=7pt, bottom=7pt, coltext=white,
  fontupper=\small\raggedright, before upper={{\color{amber}\normalsize\bfseries #2}\par\smallskip}, #1}

% Value tables: published beside Julia
\newenvironment{valuetable}[1]{%
  \par\small\tabularx{\linewidth}{#1}\toprule}{\bottomrule\endtabularx\par}
\newcommand{\ok}{{\color{agree}\faCheck}}
\newcommand{\no}{{\color{differ}\faTimes}}

% End of the document: credit and licence
\newcommand{\nmcredit}{%
  \par\vfill
  \begin{minipage}[b]{0.78\linewidth}\footnotesize\color{soft}
  Methods written and tested with Claude Code (Anthropic) in Pollis, from the published papers and their data.\par
  \textcopyright\ 2026 Antonio Saragga Seabra\enspace·\enspace Licensed under \href{https://creativecommons.org/licenses/by-sa/4.0/}{CC BY-SA 4.0}\enspace·\enspace Draft, \nmdate
  \end{minipage}\hfill
  \begin{minipage}[b]{0.18\linewidth}\raggedleft\pollislogo{9mm}\par\vspace{2pt}{\footnotesize contact@pollis.org}\end{minipage}}

\renewcommand{\nmdate}{11 November 2026}

\begin{document}

\nmhead{High-Frequency Volatility}
  {Two-Scales Realized Volatility}
  {Can a day's variance be measured from every tick when each price carries noise? Zhang, Mykland and A\"{\i}t-Sahalia
   (2005), reproduced by simulation.}

\factsheet
  {\nmcite{zma2005}{Zhang, Mykland and A\"{\i}t-Sahalia (2005)}, \emph{Journal of the American Statistical Association}:
   Table~2 (p.\,1405), five estimators on 25,000 simulated days of the model of \nmcite{heston1993}{Heston (1993)}; no
   data. Also a known volatility, and \fn{rTSCov} of the R package highfrequency 1.0.3.}
  {the two-scales estimator as the article writes it, on variances. HighFrequencyCovariance.jl
   (\nmcite{bk2022}{Baumann and Klymak 2022}) has a two-scales volatility, with the correction applied to volatilities.}
  {\menupath{Extend \gt Finance \gt High-Frequency Analysis \gt High-Frequency Volatility}: \emph{Two-Scales Volatility
   and the Noise} (\fn{tutorial-02-two-scales.ipynb}), with the Finance Toolbox}
  {43 of the 45 values of Table~2 tested: the 4 the theory gives, 3 to every printed digit and 1 off by one in the last;
   the 39 simulated, all within noise. The R package to $3.4\times10^{-6}$.}

\nmsection{The problem}
The sum of a day's squared returns estimates its integrated variance $\langle X,X\rangle_T=\int_0^T\sigma_t^2\,dt$
when prices carry no error. With noise, $Y_{t_i}=X_{t_i}+\varepsilon_{t_i}$, it grows like $2n\,E\varepsilon^2$ with
the $n$ returns: every second, it measures the noise. Five-minute returns throw most of the data away and keep a bias;
the article combines two time scales to remove it with all the data. Its Table~2 is a Monte Carlo study, so the test is
statistical: a Julia mean over simulated days, with its standard error, against the printed one.

\nmsection{The method}
Table~2 ranks five estimators: every tick, $[Y,Y]^{(\mathrm{all})}$ (fifth-best); every five minutes (fourth-best);
sparse returns at the optimal step (third-best); the average $[Y,Y]^{(\mathrm{avg})}$ of the realized variances of $K$
regular sub-grids, equation~(32) (second-best); and the two-scales estimator (first-best), equations~(55) and~(64):
\[
\widehat{\langle X,X\rangle}_T=\Bigl(1-\frac{\bar n}{n}\Bigr)^{-1}\Bigl([Y,Y]^{(\mathrm{avg})}_T-\frac{\bar n}{n}\,[Y,Y]^{(\mathrm{all})}_T\Bigr),
\qquad \bar n=\frac{n-K+1}{K},
\]
with $K=c^*n^{2/3}$ sub-grids, $c^*=\bigl(\frac{T}{12\,(E\varepsilon^2)^2}\int_0^T\sigma_t^4\,dt\bigr)^{-1/3}$, equation~(10).
The noise variance is $[Y,Y]^{(\mathrm{all})}/(2n)$, equation~(22), or $([Y,Y]^{(\mathrm{all})}-[Y,Y]^{(\mathrm{avg})})/(2(n-\bar n))$, equation~(65).
\fn{realized\_variance}, \fn{sparse\_variance}, \fn{average\_variance} and \fn{two\_scales\_variance} return the four
statistics from a day's log prices, \fn{noise\_variance} equation~(22) or, given $K$, (65), and \fn{optimal\_sparse},
\fn{optimal\_average} and \fn{optimal\_c} the tuning. \fn{HestonDesign} holds the article's design ($\mu=0.05$,
$\kappa=5$, $\alpha=0.04$, $\gamma=0.5$, $\rho=-0.5$, noise of standard deviation 0.0005, 23,400 one-second prices in a
day of $T=1/252$ years), \fn{simulate\_day} draws a day by the Euler scheme, and \fn{table2} returns Table~2 as a
$9\times5$ matrix, \fn{monte\_carlo\_table2} with its standard errors.

\nmsection{What the tests show}
A simulated value is \emph{within noise} when the gap beyond the rounding of the printed value is under 3 standard
errors of the difference of two runs, as Pollis compares random values.
\begin{valuetable}{@{}lXl@{}}
\textbf{Test} & \textbf{What is compared} & \textbf{Agreement}\\\midrule
Table 2, theory & asymptotic biases and RMSE the formulas give: 4 values & 3 to every digit, 1 in the last\\
Table 2, simulated & biases, variances, RMSEs, relative ones: 39 values & all within noise, largest 2.49 s.e.\\
Known volatility & 1\% a day, noise 0 to 20 bp, 400 days: 8 means & all within 3 s.e., largest 2.34\\
R highfrequency & \fn{rTSCov} on 4 days, 5 to 20 bp & to $3.4\times10^{-6}$\\
\end{valuetable}

\nmsub{Where they differ}
\begin{itemize}
  \item \textbf{Theory:} the fifth-best asymptotic RMSE is $\sqrt{0.0117^2+1.787\times10^{-8}}=0.0117008$, which
    rounds to 1.1701e-2; the article prints 1.1700e-2.
  \item \textbf{Simulated:} two runs differ in their digits, as they must. The five-minute small-sample variance is
    1.4099e-9 vs 1.4414e-9 (0.98 s.e.), and the first-best relative bias $-.00147$ vs $-.00045$, which carries the
    $-K/(n+1)$ of the ticks the sub-grids leave out.
  \item \textbf{Not tested:} the fifth-best relative variance and RMSE, 123,200 and 395 vs 82,502 and 340. The integrated
    variance has a Gamma law of shape 1.6, so the mean of its inverse square is infinite and they never settle
    (standard errors 79,700 and 101).
  \item \textbf{R:} highfrequency counts the $n+1$ prices where the article counts the $n$ returns, in $\bar n$ and
    $\bar n/n$: $2.1\times10^{-7}$ to $3.4\times10^{-6}$ apart, under the $10^{-5}$ tested.
\end{itemize}

\nmsection{Table 2}
Published above Julia (25,000 days); in bold the cells the theory gives, in grey the two not tested.
{\footnotesize
\begin{valuetable}{@{}Xlrrrrrr@{}}
\textbf{Row} & & \textbf{Fifth} & \textbf{Fourth} & \textbf{Third} & \textbf{Second} & \textbf{First} & \textbf{Largest $z$}\\\midrule
Small-sample bias & Published & 1.1699e-2 & 3.89e-5 & 2.18e-5 & 1.926e-5 & 2e-8 & \\
 & Julia & 1.1702e-2 & 3.91e-5 & 2.19e-5 & 1.930e-5 & -7e-8 & 1.85\\
Asymptotic bias & Published & \textbf{1.1700e-2} & \textbf{3.90e-5} & 2.20e-5 & 1.927e-5 & \textbf{0} & \\
 & Julia & \textbf{1.1700e-2} & \textbf{3.90e-5} & 2.19e-5 & 1.912e-5 & \textbf{0} & 1.63\\
Small-sample variance & Published & 1.791e-8 & 1.4414e-9 & 1.59e-9 & 9.41e-10 & 9e-11 & \\
 & Julia & 1.783e-8 & 1.4099e-9 & 1.62e-9 & 9.62e-10 & 9e-11 & 1.18\\
Asymptotic variance & Published & 1.788e-8 & 1.4409e-9 & 1.58e-9 & 9.37e-10 & 8e-11 & \\
 & Julia & 1.787e-8 & 1.4280e-9 & 1.58e-9 & 9.35e-10 & 8e-11 & 1.47\\
Small-sample RMSE & Published & 1.1699e-2 & 5.437e-5 & 4.543e-5 & 3.622e-5 & 9.4e-6 & \\
 & Julia & 1.1702e-2 & 5.423e-5 & 4.580e-5 & 3.653e-5 & 9.4e-6 & 2.49\\
Asymptotic RMSE & Published & \textbf{1.1700e-2} & 5.442e-5 & 4.546e-5 & 3.618e-5 & 8.9e-6 & \\
 & Julia & \textbf{1.1701e-2} & 5.430e-5 & 4.532e-5 & 3.606e-5 & 8.9e-6 & 0.85\\
Relative bias & Published & 182 & .61 & .18 & .15 & -.00045 & \\
 & Julia & 181 & .60 & .17 & .15 & -.00147 & 1.76\\
Relative variance & Published & {\color{soft}82502} & 1.15 & .11 & .053 & .0043 & \\
 & Julia & {\color{soft}123200} & 1.06 & .11 & .056 & .0042 & 1.48\\
Relative RMSE & Published & {\color{soft}340} & 1.24 & .37 & .28 & .065 & \\
 & Julia & {\color{soft}395} & 1.19 & .37 & .28 & .065 & 0.13\\
\end{valuetable}}

\nmsection{The ranking, reproduced}
Every tick, without correction, misses by 181 times the variance it estimates; five-minute returns by 1.19 times in
relative RMSE; the two-scales estimator by .065, as published, using every tick. The averaged estimator alone keeps a
bias of .15, which the fast scale removes.

\nmsection{Four things the article does not say}
The variance starts from its stationary Gamma law (shape 1.6, scale 0.025), which the five-minute asymptotic variance
fits (1.4280e-9 against 1.4409e-9 printed). Each day is tuned with its own quarticity and the true noise variance. The
sparse and averaged estimators cover the whole day, as the small-sample biases show. Variances are taken within
$50\times5$ bins of the days, by integrated variance and then quarticity.

\nmsection{Still open}
Figures~1 and~2 are not reproduced. Whether the optimal estimators are tuned with true or estimated values is not said;
the true ones are used here.

\nmsection{Further considerations}
\begin{itemize}
  \item \textbf{A limit.} The noise is independent of the price and from tick to tick, the prices equally spaced; a
    jump in the day enters the estimate.
  \item \textbf{An extension.} More than two scales, or dependent noise, each tested against its own article.
  \item \textbf{Upstream.} Table~2 and the known volatility are ready-made tests for a two-scales estimator on variances
    in HighFrequencyCovariance.jl.
\end{itemize}

\nmsection{References}
{\footnotesize
\nmref{zma2005}{Zhang, L., Mykland, P. A. and A\"{\i}t-Sahalia, Y. (2005). A tale of two time scales: determining integrated volatility with noisy high-frequency data. \emph{Journal of the American Statistical Association} 100(472), 1394--1411. \nmdoi{10.1198/016214505000000169}}
\nmref{bk2022}{Baumann, S. and Klymak, M. (2022). HighFrequencyCovariance: a Julia package for estimating covariance matrices using high frequency financial data. \emph{Journal of Statistical Software} 103(14). \nmdoi{10.18637/jss.v103.i14}}
\nmref{heston1993}{Heston, S. L. (1993). A closed-form solution for options with stochastic volatility with applications to bond and currency options. \emph{Review of Financial Studies} 6(2), 327--343. \nmdoi{10.1093/rfs/6.2.327}}
}

\nmcredit
\end{document}
